\chapter{Methodology}
\label{chap:methodology}

% ─────────────────────────────────────────────────────────────────────────────
\section{Overview}
\label{sec:methodology_overview}
% ─────────────────────────────────────────────────────────────────────────────

This chapter describes the design and implementation of the proposed Clinical
Multi-Agent Decisioning System (CMADS). The principal contribution of this
thesis is the multi-agent reasoning pipeline, and the remainder of this
chapter is devoted almost entirely to its design. The supporting data layer
--- a medallion lakehouse built on Synthea-generated patients --- is
summarised briefly in the Dataset section that follows (Section~\ref{sec:dataset})
so that the reader has enough context to interpret the pipeline's inputs
without being distracted from the agentic core of the system.

Figure~\ref{fig:system_architecture} shows the high-level architecture. The
data layer (top) ends at the Gold case objects; everything below that point
is the multi-agent system that this chapter is primarily concerned with.
Seven specialised agents, orchestrated by a LangGraph \texttt{StateGraph},
execute in six sequential stages to produce a final clinical decision report.

\begin{figure}[ht]
\centering
\begin{tikzpicture}[
  every node/.style={font=\scriptsize},
  data/.style={rectangle, draw, rounded corners=2pt, minimum height=10mm,
               minimum width=36mm, align=center, fill=yellow!30},
  gold/.style={rectangle, draw, rounded corners=2pt, minimum height=10mm,
               minimum width=36mm, align=center, fill=yellow!45},
  agent/.style={rectangle, draw, rounded corners=2pt, minimum height=9mm,
                minimum width=32mm, align=center, fill=green!15},
  evalbox/.style={rectangle, draw, rounded corners=2pt, minimum height=9mm,
                  minimum width=32mm, align=center, fill=red!15},
  consol/.style={rectangle, draw, rounded corners=2pt, minimum height=9mm,
                 minimum width=32mm, align=center, fill=orange!25,
                 font=\scriptsize\bfseries},
  memlayer/.style={rectangle, draw, rounded corners=4pt, minimum height=14mm,
                   minimum width=160mm, align=center, fill=violet!12,
                   font=\scriptsize\itshape},
  arr/.style={-{Latex[length=1.8mm]}, thick},
  arrgate/.style={-{Latex[length=1.6mm]}, thick, red!60!black},
  arrmem/.style={-{Latex[length=1.4mm]}, dashed, violet!50!black, thin},
  caplabel/.style={font=\scriptsize\bfseries, text=blue!50!black}
]
  % Data layer (single box) and Gold
  \node[data] (data) {\textbf{Clinical Data Layer}\\Synthea $\to$ Bronze $\to$ Silver $\to$ Gold};
  \node[gold, right=14mm of data] (goldnode) {Gold Case Objects\\(EHR + Lab JSON)};
  \draw[arr] (data) -- (goldnode);

  % Separator label "data layer"
  \node[caplabel, above=1mm of data] {\textit{Data Layer --- context only}};

  % Row of stages below (left to right) — six diagnostic stages plus
  % a seventh maintenance stage (memory consolidation, see Section
  % \ref{sec:consolidation}).
  \node[agent, below=14mm of data, xshift=6mm] (s1a) {1a.\ EHR Analyst};
  \node[agent, right=6mm of s1a] (s1b) {1b.\ Lab Interpreter};
  \node[agent, below=8mm of s1a] (s2) {2.\ Diagnostic Reasoning};
  \node[agent, right=6mm of s2] (s3) {3.\ Clinical Reviewer};
  \node[agent, below=8mm of s2] (s4) {4.\ Diagnostic Refiner};
  \node[evalbox, right=6mm of s4] (s5) {5.\ LLM Evaluator};
  \node[evalbox, below=8mm of s5, xshift=-20mm] (s6) {6.\ Treatment Planning (RAG)};
  \node[consol, below=7mm of s6] (s7) {7.\ Memory Consolidation};

  % Arrows from Gold to stage 1 agents
  \draw[arr] (goldnode.south) -- ++(0,-4mm) -| (s1a.north);
  \draw[arr] (goldnode.south) -- ++(0,-4mm) -| (s1b.north);

  % Stage 1 fan-in to Stage 2
  \draw[arr] (s1a.south) -- ++(0,-3mm) -| (s2.north);
  \draw[arr] (s1b.south) -- ++(0,-3mm) -| (s2.north);

  % Stage 2 -> 3
  \draw[arr] (s2) -- (s3);
  % Stage 3 -> 4 (goes down and left)
  \draw[arr] (s3.south) -- ++(0,-3mm) -| (s4.north);
  % Stage 4 -> 5
  \draw[arr] (s4) -- (s5);
  % Stage 5 -> 6 (DIRECT only — gated)
  \draw[arrgate] (s5.south) -- ++(0,-3mm) -| (s6.north)
    node[midway, right, font=\tiny\bfseries, xshift=2mm, yshift=3mm,
         text=red!60!black]{if DIRECT};
  % Stage 6 -> 7 (always)
  \draw[arr] (s6) -- (s7);

  \node[caplabel, below=1mm of s7] {\textit{Multi-Agent Reasoning Pipeline --- thesis focus}};

  % Multi-level memory layer (spans the full width below)
  \node[memlayer, below=7mm of s7] (mem)
    {\textbf{Multi-Level Memory Subsystem (Section~\ref{sec:multilevel_memory})}\\[1pt]
     Tier 1 Working  ·  Tier 2 Episodic  ·  Tier 3 Semantic  ·  Tier 4 Case-Based
     \quad accessed through one MemoryManager facade};

  % Dashed read/write tendrils from agents to the memory layer
  \foreach \n in {s1a, s1b, s2, s3, s4, s5, s6}
    \draw[arrmem] (\n.south) ++(0,-2mm) -- ($(mem.north)!(\n.south)!(mem.north east)$);
  \draw[arrmem, very thick, orange!70!black]
    (s7.east) -- ++(8mm, 0) |- (mem.east)
    node[midway, right, font=\tiny\bfseries, text=orange!70!black, xshift=-1mm, yshift=-1mm]
      {writes T3 + T4};
\end{tikzpicture}
\caption{High-level CMADS architecture. The data layer (top) is summarised
in Section~\ref{sec:dataset}; the seven-agent reasoning pipeline (centre)
is the thesis focus. Stage~6 runs only when Stage~5 produces a DIRECT
match. The new Stage~7 consolidation node (orange) folds the run's
outcome into the multi-level memory subsystem (Section~\ref{sec:shared_memory}),
which is read by every reasoning stage above (dashed lines).}
\label{fig:system_architecture}
\end{figure}

The chapter is organised as follows. Section~\ref{sec:dataset} briefly
describes the dataset and the data preparation pipeline that produces the
Gold case objects. Section~\ref{sec:mas_framework} introduces the multi-agent
framework and the LangGraph orchestrator. Section~\ref{sec:shared_memory}
details the shared memory design. Section~\ref{sec:agent_blueprint} presents
the generic agent blueprint. Section~\ref{sec:agents} describes each of the
seven specialised agents. Section~\ref{sec:rag} covers the
Retrieval-Augmented Generation layer for treatment planning.
Section~\ref{sec:llm_config} documents the LLM provider abstraction.
Section~\ref{sec:evaluation} defines the evaluation framework.
Section~\ref{sec:env} lists the hardware and software environment, and
Section~\ref{sec:pipeline_e2e} summarises the end-to-end execution flow.

% ─────────────────────────────────────────────────────────────────────────────
\section{Dataset}
\label{sec:dataset}
% ─────────────────────────────────────────────────────────────────────────────

The patient records consumed by the multi-agent pipeline are produced by
\textit{Synthea}, an open-source rule-based patient simulator developed by
the MITRE Corporation. Synthea is a deterministic, state-machine-based
generator informed by epidemiological data from the Centers for Disease
Control and the National Institutes of Health, and it emits patient records
in the FHIR R4 standard. Two properties motivated this choice: Synthea
guarantees ground truth for every condition, medication, and observation it
emits, which is essential for the LLM-as-Judge evaluation described in
Section~\ref{sec:evaluation}; and being fully synthetic, it has no ethical
or data-governance constraints, enabling open reproducibility.

A cohort of 10{,}000 patients was generated, from which 270 clinically
diverse cases covering eight target diseases --- ischemic heart disease,
congestive heart failure, essential hypertension, diabetes mellitus type 2,
chronic kidney disease stages 2 and 3, end-stage renal disease, and
metabolic syndrome X --- were selected by a deterministic target-encounter
algorithm that retains only patients with at least three prior encounters
and five laboratory observations before the date of diagnosis.

The raw FHIR output passes through a four-stage medallion lakehouse before
reaching the agents. Bronze stores the raw Parquet-encoded FHIR resources;
Silver transforms them into the Observational Medical Outcomes Partnership
Common Data Model (OMOP CDM v5.4) using dbt-core and DuckDB;
Silver\textsuperscript{+} adds pre-computed derived features such as
longitudinal laboratory trends and estimated glomerular filtration rate; and
the Gold layer assembles the OMOP relational data into two JSON artefacts
per patient: an \texttt{ehr\_case.json} containing demographics, active
conditions, and history, and a \texttt{lab\_case.json} containing laboratory
time-series with flagged critical values. A third \texttt{ground\_truth.json}
file stores the Synthea target condition and is never exposed to any
reasoning agent. The Gold case objects are the sole input to the multi-agent
pipeline and are the boundary at which the data layer ends and the agentic
system begins.

% ─────────────────────────────────────────────────────────────────────────────
\section{Multi-Agent Framework Architecture}
\label{sec:mas_framework}
% ─────────────────────────────────────────────────────────────────────────────

The multi-agent pipeline is orchestrated by \textit{LangGraph}, a graph-based
workflow engine that models agent execution as a typed state machine. Each
agent is a node in a \texttt{StateGraph}, edges encode execution order and
conditional routing, and a shared \texttt{PipelineState} dictionary flows
through the graph carrying the cumulative output of every previously executed
agent. The choice of LangGraph over a hand-written orchestrator was motivated
by three properties: first-class support for parallel fan-out and fan-in,
which Stage~1 requires; a checkpointing hook that persists the full execution
state after every agent and supports recovery from partial failures; and
reducer-based state semantics that prevent one agent from silently
overwriting another agent's output, directly addressing the concurrency
corruption risks identified by Chen et al.\ in
MedSentry~\cite{Chen2025MedSentry}.

Figure~\ref{fig:agent_pipeline} shows the complete execution graph. The
pipeline is decomposed into six sequential stages, with Stage~1 running two
agents in parallel and Stage~2 implementing an adaptive self-critique loop
bounded at three rounds.

\begin{figure}[ht]
\centering
\begin{tikzpicture}[
  every node/.style={font=\scriptsize},
  agent/.style={rectangle, draw, rounded corners=2pt, minimum height=10mm,
                minimum width=46mm, align=center, fill=green!15},
  consol/.style={rectangle, draw, rounded corners=2pt, minimum height=10mm,
                 minimum width=46mm, align=center, fill=orange!25,
                 font=\scriptsize\bfseries},
  stage/.style={rectangle, draw=blue!40, rounded corners=2pt, minimum height=7mm,
                minimum width=18mm, align=center, fill=blue!10, font=\scriptsize\bfseries},
  memside/.style={rectangle, draw=violet!50, rounded corners=3pt,
                  minimum height=86mm, minimum width=26mm, align=center,
                  fill=violet!10, font=\scriptsize, inner sep=2pt},
  tiertag/.style={font=\scriptsize\bfseries, text=violet!60!black},
  arr/.style={-{Latex[length=1.6mm]}, thick},
  loopback/.style={-{Latex[length=1.6mm]}, thick, dashed, blue!60!black},
  memarrow/.style={-{Latex[length=1.4mm]}, thin, dashed, violet!60!black}
]
  % Starting node
  \node[agent, fill=yellow!40] (gold) {Gold Case Object (EHR + Lab)};

  % Stage 1 (parallel)
  \node[agent, below left=8mm and -14mm of gold] (ehr) {1a.\ EHR Analyst};
  \node[agent, below right=8mm and -14mm of gold] (lab) {1b.\ Lab Interpreter};

  % Stage 2
  \node[agent, below=22mm of gold] (diag) {2.\ Diagnostic Reasoning\\(adaptive loop / 3-phase split)};

  % Stage 3
  \node[agent, below=8mm of diag] (rev) {3.\ Clinical Reviewer (adversarial)};

  % Stage 4
  \node[agent, below=8mm of rev] (ref) {4.\ Diagnostic Refiner};

  % Stage 5
  \node[agent, below=8mm of ref, fill=red!15] (evalnode) {5.\ LLM Evaluator vs.\ Ground Truth};

  % Stage 6
  \node[agent, below=8mm of evalnode, fill=red!15] (tx) {6.\ Treatment Planning (RAG)};

  % Stage 7 (NEW — memory consolidation)
  \node[consol, below=8mm of tx] (mc) {7.\ Memory Consolidation\\(non-LLM maintenance)};

  % Edges
  \draw[arr] (gold.south) -- ++(0,-3mm) -| (ehr.north);
  \draw[arr] (gold.south) -- ++(0,-3mm) -| (lab.north);
  \draw[arr] (ehr.south) |- ($(diag.north)+(-14mm,3mm)$) -- ++(0,-3mm);
  \draw[arr] (lab.south) |- ($(diag.north)+(14mm,3mm)$) -- ++(0,-3mm);
  \draw[arr] (diag) -- (rev);
  \draw[arr] (rev) -- (ref);
  \draw[arr] (ref) -- (evalnode);
  \draw[arr] (evalnode) -- node[right, font=\tiny, xshift=1mm]{if DIRECT} (tx);
  \draw[arr] (tx) -- (mc);

  % Adaptive / 3-phase loopback at Stage 2
  \draw[loopback] ($(diag.east)+(0,2mm)$)
    .. controls +(right:14mm) and +(right:14mm) ..
    ($(diag.east)+(0,-2mm)$)
    node[midway, right=3mm, align=left, font=\tiny]
        {critique\\($\leq 3$ rounds)};

  % Stage labels (left column)
  \node[stage, left=10mm of ehr, xshift=-4mm] (s1) {Stage 1};
  \node[stage, left=10mm of diag, xshift=-4mm] (s2) {Stage 2};
  \node[stage, left=10mm of rev, xshift=-4mm] (s3) {Stage 3};
  \node[stage, left=10mm of ref, xshift=-4mm] (s4) {Stage 4};
  \node[stage, left=10mm of evalnode, xshift=-4mm] (s5) {Stage 5};
  \node[stage, left=10mm of tx, xshift=-4mm] (s6) {Stage 6};
  \node[stage, left=10mm of mc, xshift=-4mm, fill=orange!25] (s7) {Stage 7};

  % Memory column on the right — shared by all stages
  \node[memside, right=14mm of diag] (memcol)
    {\textbf{Multi-Level Memory}\\[2pt]
     \textcolor{orange!70!black}{\textbf{T1 Working}}\\
     used by Stage 2\\
     critique loop\\[2pt]
     \textcolor{blue!60!black}{\textbf{T2 Episodic}}\\
     all stages\\
     append timeline\\[2pt]
     \textcolor{teal!60!black}{\textbf{T3 Semantic}}\\
     Refiner / Tx\\
     read priors\\[2pt]
     \textcolor{violet!60!black}{\textbf{T4 Case-based}}\\
     Diag reads\\
     similar past\\[2pt]
     \textit{Stage 7 writes T3 \& T4}};

  % Per-stage dashed connectors to the memory column
  \foreach \n in {ehr, lab, diag, rev, ref, evalnode, tx}
    \draw[memarrow] (\n.east) -- ($(memcol.west)!(\n.east)!(memcol.south west)$);
  \draw[memarrow, very thick, orange!70!black]
    (mc.east) -- ($(memcol.west)!(mc.east)!(memcol.south west)$);

\end{tikzpicture}
\caption{Multi-agent execution graph with the Stage~7 consolidation node
and the multi-level memory column (right). Stage~1 fans out the EHR
Analyst and Lab Interpreter in parallel; Stage~2 contains an adaptive
self-critique loop bounded at three rounds (and a 3-phase
with-prior / clean-room / synthesise variant when Tier-4 priors are
available, Section~\ref{sec:diagnostic_reasoning}); Stage~6 runs only
when Stage~5 produces a DIRECT match. Dashed lines indicate read/write
access to the memory subsystem; the orange link from Stage~7 indicates
the per-run consolidation writes into Tier~3 and Tier~4.}
\label{fig:agent_pipeline}
\end{figure}

The pipeline comprises six diagnostic stages plus a seventh maintenance
stage that consolidates each run's outcome into the cross-session
memory tiers (Section~\ref{sec:consolidation}). The seven stages are
summarised below.
\begin{itemize}
  \item \textbf{Stage 1 --- Parallel context extraction:} the EHR Analyst and
  Lab Interpreter run concurrently, each consuming its Gold case file and
  producing a structured clinical summary.
  \item \textbf{Stage 2 --- Diagnostic reasoning:} an adaptive loop generates
  a ranked differential of at least three candidates, self-critiques it, and
  re-ranks for up to three rounds or until a confidence threshold of 75\% is
  reached.
  \item \textbf{Stage 3 --- Clinical review:} an adversarial reviewer agent
  independently re-examines the evidence and issues per-diagnosis verification
  flags.
  \item \textbf{Stage 4 --- Diagnostic refinement:} the outputs of Stages~2
  and 3 are merged into a single final differential.
  \item \textbf{Stage 5 --- Evaluation:} the final differential is compared
  against the Synthea ground truth by an LLM-as-Judge evaluator that assigns
  DIRECT, INDIRECT, or MISS classifications.
  \item \textbf{Stage 6 --- Treatment planning:} for DIRECT matches only, a
  treatment agent retrieves relevant NICE guideline passages from a Qdrant
  vector store and assembles a guideline-grounded treatment plan.
  \item \textbf{Stage 7 --- Memory consolidation:} a non-LLM maintenance node
  folds the run's outcome into the cross-session semantic store (Tier~3) and
  the case-based vector store (Tier~4); see Section~\ref{sec:consolidation}.
\end{itemize}

% ─────────────────────────────────────────────────────────────────────────────
\section{Shared Memory Design}
\label{sec:shared_memory}
% ─────────────────────────────────────────────────────────────────────────────

Inter-agent communication is mediated through a typed shared memory
implemented as a \texttt{TypedDict} named \texttt{PipelineState}, layered
beneath a four-tier multi-level memory subsystem inspired by the
\emph{Cognitive Architectures for Language Agents} (CoALA) framework
\cite{Sumers2024CoALA}. The base layer provides deterministic merging
of partial updates within a single patient run; the four tiers above it
provide intra-agent scratch space, an inter-agent reasoning trail, a
cross-session disease-level prior, and a vector-based case-similarity
recall over previously processed patients.

% ─────────────────────────────────────────────────────────────────────────────
\subsection{Base layer: \texttt{PipelineState}}
\label{sec:pipeline_state}
% ─────────────────────────────────────────────────────────────────────────────

Each field of \texttt{PipelineState} is a separate namespace annotated
with an explicit reducer function that defines how concurrent partial
updates are merged. This makes concurrent-write conflicts impossible by
construction: when two parallel agents write to the same namespace, the
reducer resolves the merge deterministically rather than producing an
undefined last-writer-wins state.

Seven namespaces partition the shared memory; the first five form the
historical core, and two are added for the episodic memory tier
introduced in Section~\ref{sec:tier2}:

\begin{itemize}
  \item \texttt{patient\_context} --- the Gold case object, written once by
  the orchestrator and read by all agents.
  \item \texttt{agent\_outputs} --- a dictionary with one slot per agent,
  using a write-own / read-upstream access pattern. The merge reducer
  ensures that parallel Stage-1 agents never overwrite each other.
  \item \texttt{execution\_trace} --- an append-only invocation log
  recording the status, duration, and error of every agent run.
  \item \texttt{scratchpad} --- per-agent dictionaries; in the present
  design this slot is repurposed as the working-memory tier
  (Section~\ref{sec:tier1}).
  \item \texttt{conflicts} --- an append-only list reserved for an
  orchestrator-level diff engine; not currently written by any agent and
  retained for forward compatibility.
  \item \texttt{session\_memory} --- append-only list of typed
  \texttt{SessionEvent} dictionaries that constitute the episodic-memory
  timeline for the current run.
  \item \texttt{session\_summary} --- per-agent one-line digest with a
  merge reducer, intended as a lightweight scoreboard for downstream
  routing logic.
\end{itemize}

% ─────────────────────────────────────────────────────────────────────────────
\subsection{Multi-level memory subsystem}
\label{sec:multilevel_memory}
% ─────────────────────────────────────────────────────────────────────────────

Restricting agent communication to the namespaces above exposes a
substantive limitation: downstream agents see only the final structured
output of upstream agents and have no visibility into the reasoning that
produced it, and no information of any kind flows from one patient run
to the next. The four-tier memory subsystem layered on top of
\texttt{PipelineState} is designed to address both problems while
preserving the existing reducer guarantees. The four tiers are scoped
from narrowest to broadest in lifespan and reuse, and are accessed
through a single \texttt{MemoryManager} facade so that agent code
follows the same call pattern regardless of which tier it touches. The
entire subsystem is gated by a \texttt{MEMORY\_ENABLED} environment
flag, which makes before/after experimentation a clean toggle.
Figure~\ref{fig:multilevel_memory} sketches the resulting architecture.

\begin{figure}[ht]
\centering
\begin{tikzpicture}[
  every node/.style={font=\scriptsize},
  agent/.style={rectangle, draw, rounded corners=2pt, minimum height=8mm,
                minimum width=22mm, align=center, fill=green!15,
                font=\scriptsize\bfseries},
  newagent/.style={rectangle, draw, rounded corners=2pt, minimum height=8mm,
                   minimum width=24mm, align=center, fill=orange!25,
                   font=\scriptsize\bfseries},
  facade/.style={rectangle, draw, rounded corners=4pt, minimum height=11mm,
                 minimum width=148mm, align=center,
                 fill=violet!22, font=\scriptsize\bfseries},
  t1/.style={rectangle, draw, rounded corners=3pt, minimum height=18mm,
             minimum width=33mm, align=center, fill=orange!18,
             font=\scriptsize},
  t2/.style={rectangle, draw, rounded corners=3pt, minimum height=18mm,
             minimum width=33mm, align=center, fill=blue!12,
             font=\scriptsize},
  t3/.style={rectangle, draw, rounded corners=3pt, minimum height=18mm,
             minimum width=33mm, align=center, fill=teal!18,
             font=\scriptsize},
  t4/.style={rectangle, draw, rounded corners=3pt, minimum height=18mm,
             minimum width=33mm, align=center, fill=violet!18,
             font=\scriptsize},
  ext/.style={rectangle, draw, dashed, rounded corners=3pt,
              minimum height=10mm, minimum width=148mm,
              align=center, fill=gray!8, font=\scriptsize\itshape},
  scope/.style={font=\tiny\itshape, text=gray!50!black},
  arrdash/.style={-{Latex[length=1.4mm]}, dashed, gray!70, thin},
  arr/.style={-{Latex[length=1.6mm]}, semithick}
]
  % ── Agent row (top) ──
  \node[agent] (a1) {EHR};
  \node[agent, right=2mm of a1] (a2) {Lab};
  \node[agent, right=2mm of a2] (a3) {Diagnostic};
  \node[agent, right=2mm of a3] (a4) {Reviewer};
  \node[agent, right=2mm of a4] (a5) {Refiner};
  \node[agent, right=2mm of a5] (a6) {Evaluator};
  \node[agent, right=2mm of a6] (a7) {Treatment};
  \node[newagent, right=2mm of a7] (a8) {Consolidator};

  % ── MemoryManager facade ──
  \node[facade, below=4mm of a4, xshift=12mm] (facade)
    {\texttt{MemoryManager} facade
     \quad\textnormal{\scriptsize\itshape one entry point: \texttt{mm.working},
     \texttt{mm.episodic}, \texttt{mm.semantic}, \texttt{mm.case\_based}}};

  % Dashed read/write lines from each agent to the facade
  \foreach \src in {a1, a2, a3, a4, a5, a6, a7, a8}
    \draw[arrdash] (\src.south) -- ($(facade.north)!(\src.south)!(facade.north east)$);

  % ── Four tier cards ──
  \node[t1, below=12mm of facade.south west, xshift=18mm] (t1)
    {\textbf{Tier 1}\\\textbf{WORKING}\\[1pt]
     \scriptsize per-agent\\
     \scriptsize invocation-local\\[2pt]
     \scriptsize \texttt{state.scratchpad}};
  \node[t2, right=4mm of t1] (t2)
    {\textbf{Tier 2}\\\textbf{EPISODIC}\\[1pt]
     \scriptsize current run\\
     \scriptsize in-state, on disk\\[2pt]
     \scriptsize \texttt{session\_memory}};
  \node[t3, right=4mm of t2] (t3)
    {\textbf{Tier 3}\\\textbf{SEMANTIC}\\[1pt]
     \scriptsize cross-session\\
     \scriptsize disk · per-disease\\[2pt]
     \scriptsize \texttt{semantic\_memory.json}};
  \node[t4, right=4mm of t3] (t4)
    {\textbf{Tier 4}\\\textbf{CASE-BASED}\\[1pt]
     \scriptsize cross-session\\
     \scriptsize Qdrant · BioLORD\\[2pt]
     \scriptsize \texttt{patient\_cases}};

  % Arrows from facade down to each tier
  \draw[arr] (facade.south) -- ($(t1.north)+(0,1mm)$);
  \draw[arr] (facade.south) -- ($(t2.north)+(0,1mm)$);
  \draw[arr] (facade.south) -- ($(t3.north)+(0,1mm)$);
  \draw[arr] (facade.south) -- ($(t4.north)+(0,1mm)$);

  % ── Scope axis label below tier cards ──
  \node[scope, below=2mm of t1] (lt1) {one LLM call};
  \node[scope, below=2mm of t2] (lt2) {one patient run};
  \node[scope, below=2mm of t3] (lt3) {disease-level aggregate};
  \node[scope, below=2mm of t4] (lt4) {case-similarity recall};
  \draw[->, gray!60, thin] ($(lt1.south west)+(-3mm,-1mm)$) --
        ($(lt4.south east)+(3mm,-1mm)$)
        node[midway, below=0mm, font=\tiny\itshape, text=gray!70!black]
            {scope $\rightarrow$ longer lifespan, broader reuse};

  % ── External knowledge base (NICE) shown but excluded from tiers ──
  \node[ext, below=12mm of t1.south, xshift=63mm] (nice)
    {\textbf{External knowledge base — not a memory tier} \quad
     NICE clinical guidelines (Qdrant \texttt{nice\_guidelines}),
     used by the Treatment agent only};

\end{tikzpicture}
\caption{Four-tier multi-level memory subsystem layered on the
LangGraph state. All seven specialised agents (green) and the new
Stage-7 memory consolidation node (orange) reach memory through a
single \texttt{MemoryManager} facade (purple) that exposes one method
per tier. Tiers are scoped from narrowest (per-LLM-call working scratch)
to broadest (cross-session vector recall over previously processed
patients). NICE clinical guidelines remain available but live outside
the memory hierarchy as static external reference knowledge.}
\label{fig:multilevel_memory}
\end{figure}

% ── Tier 1 ───────────────────────────────────────────────────────────────────
\paragraph{Tier 1 --- Working memory.}
\label{sec:tier1}
Per-agent scratch space, scoped to a single agent invocation. Lives in
the previously unused \texttt{scratchpad[agent\_id]} slot of
\texttt{PipelineState}. The Diagnostic Reasoning agent
(Section~\ref{sec:diagnostic_reasoning}) is the principal user: across
its adaptive critique loop it persists the per-round confidence value
and a rolling list of critique snippets so that subsequent calls within
the same invocation can reason about earlier rounds without reparsing
them from free-text concatenation. Working memory is cleared at the
start of every invocation and is never visible to other agents---that
role belongs to Tier~2.

% ── Tier 2 ───────────────────────────────────────────────────────────────────
\paragraph{Tier 2 --- Episodic memory.}
\label{sec:tier2}
A typed timeline of significant events across one patient run, written
to the new \texttt{session\_memory} channel through an append reducer.
Each entry is a Pydantic-validated \texttt{SessionEvent} record with an
\texttt{event\_type} drawn from a fixed vocabulary (\texttt{critique},
\texttt{confidence\_check}, \texttt{decision}, \texttt{hypothesis},
\texttt{evidence\_link}, and lifecycle markers), an \texttt{agent\_id},
a one-line summary, and a structured payload. The Clinical Reviewer
consumes the critique-trail subset of these events when challenging the
Diagnostic agent's differential, and the Diagnostic Refiner consumes
the full timeline before producing the final differential.
Lifecycle events plus a per-agent one-line digest are also written to
\texttt{session\_summary} as a cheap status channel; the timeline itself
is persisted to disk per patient as \texttt{session\_memory.json}
alongside the existing \texttt{execution\_trace.json}.

% ── Tier 3 ───────────────────────────────────────────────────────────────────
\paragraph{Tier 3 --- Semantic memory.}
\label{sec:tier3}
Per-disease aggregate statistics maintained across patient runs and
persisted to a single JSON file
(\texttt{data/gold/memory/semantic\_memory.json}). For each disease
encountered, the store accumulates the count of \textsc{direct},
\textsc{indirect}, and \textsc{miss} evaluations, the rank-1 frequency
when found, the running mean of the top differential probability, and a
capped set of observed evidence patterns. Updates are atomic
(temporary-file write followed by rename) and idempotent on repeated
re-indexing of the same patient. Recall is direct lookup by disease
name (case-insensitive); the file scale and update cadence do not
warrant an embedding-based index. The Diagnostic Refiner and the
Treatment Planning agent read the relevant entries as Bayesian-style
priors over the diseases that appear in the current differential.

% ── Tier 4 ───────────────────────────────────────────────────────────────────
\paragraph{Tier 4 --- Case-based memory.}
\label{sec:tier4}
A vector store of past patient cases, providing case-similarity recall
that the disk-backed Tier~3 cannot. After every completed run, a short
clinical-feature text---demographics, active conditions, active
medications, and key abnormal labs---is embedded with the BioLORD-2023
biomedical sentence encoder \cite{Remy2023BioLORD} and upserted to a
dedicated Qdrant collection (\texttt{patient\_cases}) keyed by a stable
hash of the patient identifier. The payload stored alongside the
embedding contains the matched diagnosis, match type, rank, primary
confidence, and a small set of evidence-pattern strings.
At the start of a new patient's reasoning, the Diagnostic agent issues
a top-$K$ similarity query (cosine distance, $K=5$ by default) and
injects the retrieved cases into its prompt as priors of the form
``patients with this profile previously landed on diagnosis $d$ with
match $t$''. The wrapper degrades gracefully: if the Qdrant cluster is
unreachable or the embedding model is not loaded, recall returns an
empty list rather than aborting the run.

The NICE clinical guidelines used by the Treatment Planning agent
(Section~\ref{sec:treatment}) live in a separate Qdrant collection
(\texttt{nice\_guidelines}) and are not part of the memory hierarchy:
they are static external reference knowledge, not memory the system
accumulates from experience. The same \texttt{MemoryManager} exposes
them through a \texttt{knowledge\_base} attribute for symmetry, but
they are excluded from the access table below.

% ─────────────────────────────────────────────────────────────────────────────
\subsection{Memory consolidation node}
\label{sec:consolidation}
% ─────────────────────────────────────────────────────────────────────────────

Tier~3 and Tier~4 are populated by a dedicated maintenance node placed
as Stage~7 of the LangGraph pipeline, executed unconditionally after
the Treatment Planning agent. The node reads the run's final
diagnosis, the evaluator's verdict, and the patient context, then
performs two writes: a per-disease consolidation into Tier~3 and a
patient-case upsert into Tier~4. The two writes are independent; if
one fails (for example, when the Qdrant cluster is unreachable), the
other still proceeds and the trace records the partial outcome. This
isolates the memory subsystem from the diagnostic pipeline so that
infrastructure failures cannot corrupt patient outputs already written
by Stages~1--6.

% ─────────────────────────────────────────────────────────────────────────────
\subsection{Access pattern}
\label{sec:memory_access}
% ─────────────────────────────────────────────────────────────────────────────

Table~\ref{tab:memory_access} summarises the read/write access pattern
of each agent across the four memory tiers. The historical invariant
is preserved: every agent writes only to its own slot in
\texttt{agent\_outputs} and to the append-only Tier-2 timeline,
eliminating concurrent-write conflicts by construction. Cross-session
writes (Tiers~3 and 4) are localised to the consolidation node, which
runs after all diagnostic stages have completed.

\begin{table}[ht]
\centering
\caption{Multi-level memory access pattern. R = reads, W = writes (to
own slot for \textsc{tier 2}, append for \textsc{tier 3}/\textsc{tier 4}).
The Treatment agent additionally reads the external NICE-guideline
knowledge base, which is not classified as memory.}
\label{tab:memory_access}
\resizebox{\textwidth}{!}{%
\begin{tabular}{@{}lcccc@{}}
\toprule
\textbf{Agent} & \textbf{Tier 1 Working} & \textbf{Tier 2 Episodic} & \textbf{Tier 3 Semantic} & \textbf{Tier 4 Case-based} \\
\midrule
EHR Analyst              & --   & W            & --  & -- \\
Lab Interpreter          & --   & W            & --  & -- \\
Diagnostic Reasoning     & RW   & R~/~W        & --  & R \\
Clinical Reviewer        & --   & R~(critique) & --  & -- \\
Diagnostic Refiner       & --   & R~/~W        & R   & -- \\
LLM Evaluator            & --   & W            & --  & -- \\
Treatment Planning       & --   & W            & R   & -- \\
Memory Consolidation~(S7)& --   & W            & W   & W \\
\bottomrule
\end{tabular}%
}
\end{table}

The implementation lives under \texttt{src/memory/} with one module per
tier (\texttt{working\_memory.py}, \texttt{episodic\_memory.py},
\texttt{semantic\_memory.py}, \texttt{case\_based\_memory.py}) and a
unified \texttt{MemoryManager} facade in \texttt{manager.py}. The
consolidation node is implemented as
\texttt{src/agents/memory\_consolidator.py} and wired into the graph
between the Treatment Planning agent and the \texttt{END} node.

% ─────────────────────────────────────────────────────────────────────────────
\section{Agent Blueprint}
\label{sec:agent_blueprint}
% ─────────────────────────────────────────────────────────────────────────────

Every specialised agent in the pipeline shares a common internal architecture,
illustrated in Figure~\ref{fig:agent_blueprint}. The blueprint is implemented
as a Python base class and specialised by subclasses that supply
agent-specific prompts, schemas, and post-processing logic. The uniformity of
this blueprint is a deliberate design decision: it ensures that every agent
exposes the same failure modes, tracing hooks, and retry semantics, which
substantially simplifies debugging and observability.

\begin{figure}[ht]
\centering
\begin{tikzpicture}[
  every node/.style={font=\scriptsize},
  gate/.style={rectangle, draw, rounded corners=2pt, minimum height=11mm,
               minimum width=23mm, align=center, fill=orange!20},
  mod/.style={rectangle, draw, rounded corners=2pt, minimum height=11mm,
              minimum width=23mm, align=center, fill=blue!10},
  llm/.style={rectangle, draw, rounded corners=2pt, minimum height=11mm,
              minimum width=23mm, align=center, fill=red!15},
  facade/.style={rectangle, draw, rounded corners=3pt, minimum height=8mm,
                 minimum width=119mm, align=center, fill=violet!20,
                 font=\scriptsize\bfseries},
  t1/.style={rectangle, draw, rounded corners=2pt, minimum height=10mm,
             minimum width=27mm, align=center, fill=orange!18},
  t2/.style={rectangle, draw, rounded corners=2pt, minimum height=10mm,
             minimum width=27mm, align=center, fill=blue!12},
  t3/.style={rectangle, draw, rounded corners=2pt, minimum height=10mm,
             minimum width=27mm, align=center, fill=teal!18},
  t4/.style={rectangle, draw, rounded corners=2pt, minimum height=10mm,
             minimum width=27mm, align=center, fill=violet!18},
  arr/.style={-{Latex[length=1.8mm]}, thick},
  darr/.style={-{Latex[length=1.6mm]}, thick, dashed, gray!80},
  memarr/.style={-{Latex[length=1.4mm]}, dashed, violet!50!black, thin}
]
  % Pipeline row
  \node[gate] (in) {Input Gate\\\textit{Pydantic validate}};
  \node[mod, right=3mm of in] (prompt) {Prompt Assembler\\\textit{YAML template}};
  \node[llm, right=3mm of prompt] (llm) {LLM Invocation\\\textit{retry + timeout}};
  \node[mod, right=3mm of llm] (parse) {Output Parser\\\textit{JSON $\to$ Pydantic}};
  \node[gate, right=3mm of parse] (out) {Output Gate\\\textit{clinical checks}};

  \draw[arr] (in) -- (prompt);
  \draw[arr] (prompt) -- (llm);
  \draw[arr] (llm) -- (parse);
  \draw[arr] (parse) -- (out);

  % MemoryManager facade row
  \node[facade, below=10mm of llm] (facade)
    {\texttt{MemoryManager} facade};

  % Four-tier row beneath the facade
  \node[t1, below=6mm of facade.south west, anchor=north west] (t1)
    {\textbf{Tier 1}\\Working};
  \node[t2, right=3mm of t1] (t2)
    {\textbf{Tier 2}\\Episodic};
  \node[t3, right=3mm of t2] (t3)
    {\textbf{Tier 3}\\Semantic};
  \node[t4, right=3mm of t3] (t4)
    {\textbf{Tier 4}\\Case-based};

  % Read/write links from gates to facade
  \draw[darr] (in.south) -- (facade.north -| in.south)
    node[midway, right=1mm, font=\tiny]{read};
  \draw[darr] (prompt.south) -- (facade.north -| prompt.south)
    node[midway, right=1mm, font=\tiny]{read priors};
  \draw[darr] (out.south) -- (facade.north -| out.south)
    node[midway, right=1mm, font=\tiny]{write};

  % Facade fan-out to each tier
  \draw[memarr] (facade.south) -- (t1.north);
  \draw[memarr] (facade.south) -- (t2.north);
  \draw[memarr] (facade.south) -- (t3.north);
  \draw[memarr] (facade.south) -- (t4.north);

  % Input/output labels
  \node[above=2mm of in, font=\scriptsize] {\textit{from upstream}};
  \node[above=2mm of out, font=\scriptsize] {\textit{to downstream}};
\end{tikzpicture}
\caption{Generic agent blueprint. Every agent executes Input~Gate~$\to$~Prompt
Assembler~$\to$~LLM~Invocation~$\to$~Output~Parser~$\to$~Output~Gate.
Upstream namespaces are read through the Input Gate; Tier-3 and Tier-4
priors are pulled through the Prompt Assembler when relevant
(Section~\ref{sec:multilevel_memory}); the final validated output and any
episodic events are written through the Output Gate. All four memory
tiers are reached through a single \texttt{MemoryManager} facade so the
agent code is identical regardless of which tier is touched.}
\label{fig:agent_blueprint}
\end{figure}

The five stages of the blueprint are as follows.
\begin{enumerate}
  \item The \textbf{Input Gate} validates that every required upstream
  namespace is present and conforms to its Pydantic schema. If validation
  fails, the agent records a structured error in the \texttt{execution\_trace}
  and terminates gracefully, allowing downstream agents to continue with
  partial results.
  \item The \textbf{Prompt Assembler} constructs the LLM prompt from a
  versioned YAML template, substituting the validated upstream outputs.
  Keeping prompts in YAML rather than hard-coded Python strings makes prompt
  engineering a configuration activity, supporting rapid iteration and
  version control.
  \item The \textbf{LLM Invocation} layer dispatches the prompt to the
  configured language model through the provider-agnostic adapter
  (Section~\ref{sec:llm_config}). Every call is wrapped in an exponential
  back-off retry policy (up to three attempts) and a per-call timeout.
  \item The \textbf{Output Parser} converts the raw LLM response to a
  Pydantic v2 structured object, applying JSON repair heuristics for common
  formatting failures.
  \item The \textbf{Output Gate} applies clinical-plausibility checks before
  releasing the object to shared memory. For example, the Diagnostic
  Reasoning Output Gate rejects differentials that contain fewer than three
  ranked diagnoses, and the Treatment Planning Output Gate rejects plans that
  reference medications outside the retrieved NICE guideline corpus.
\end{enumerate}

% ─────────────────────────────────────────────────────────────────────────────
\section{Agent Specifications}
\label{sec:agents}
% ─────────────────────────────────────────────────────────────────────────────

This section describes each of the seven specialised agents in detail. For
every agent the upstream inputs, downstream outputs, and distinctive
implementation features are documented. Table~\ref{tab:agent_roles}
summarises the agents and their shared-memory dependencies.

\begin{table}[ht]
\centering
\caption{The seven specialised agents and their shared-memory dependencies.}
\label{tab:agent_roles}
\small
\begin{tabular}{@{}p{30mm}p{40mm}p{40mm}p{30mm}@{}}
\toprule
\textbf{Agent} & \textbf{Reads} & \textbf{Writes} & \textbf{Distinctive} \\
\midrule
EHR Analyst          & ehr\_case                              & ehr\_analyst       & 3-call sequence \\
Lab Interpreter      & lab\_case                              & lab\_interpreter   & severity ranking \\
Diagnostic Reasoning & EHR + Lab outputs                      & diagnostic\_reasoning & adaptive loop \\
Clinical Reviewer    & Stages 1 \& 2 outputs                  & clinical\_reviewer & adversarial \\
Diagnostic Refiner   & Diagnostic + Reviewer                  & final\_diagnosis   & merge logic \\
LLM Evaluator        & Final diagnosis + ground truth         & evaluation         & separate LLM \\
Treatment Planning   & Final diagnosis + evaluation           & treatment\_planning & RAG + Qdrant \\
\bottomrule
\end{tabular}
\end{table}

\subsection{EHR Analyst}
\label{sec:agent_ehr}
The EHR Analyst consumes the \texttt{ehr\_case.json} Gold object and produces
a structured clinical summary covering demographics, active conditions,
medication list, relevant history, and prior encounters. The agent is
prompted to think step-by-step through the EHR and to cite the specific
record that supports every structured claim, which permits downstream agents
and human reviewers to verify provenance. Internally the agent runs three
sequential LLM calls within its blueprint (Figure~\ref{fig:agent_ehr_flow}):
an unstructured \textit{analysis} call, a \textit{structure} call that
extracts JSON fields, and a \textit{review} call that self-checks the
structured output for internal consistency.

\begin{figure}[ht]
\centering
\begin{tikzpicture}[
  every node/.style={font=\scriptsize},
  inp/.style={rectangle, draw, rounded corners=2pt, minimum height=10mm,
              minimum width=24mm, align=center, fill=yellow!30},
  llm/.style={rectangle, draw, rounded corners=2pt, minimum height=10mm,
              minimum width=24mm, align=center, fill=blue!12},
  outbox/.style={rectangle, draw, rounded corners=2pt, minimum height=10mm,
              minimum width=28mm, align=center, fill=green!18},
  arr/.style={-{Latex[length=1.6mm]}, semithick}
]
  \node[inp] (in) {\texttt{ehr\_case.json}\\(Gold object)};
  \node[llm, right=4mm of in] (a)
    {Call 1: \textbf{analyse}\\\scriptsize free-text reasoning};
  \node[llm, right=4mm of a] (s)
    {Call 2: \textbf{structure}\\\scriptsize extract JSON fields};
  \node[llm, right=4mm of s] (r)
    {Call 3: \textbf{review}\\\scriptsize self-check consistency};
  \node[outbox, right=4mm of r] (o)
    {\texttt{ehr\_analyst}\\\scriptsize structured summary\\\scriptsize + provenance};
  \draw[arr] (in) -- (a); \draw[arr] (a) -- (s);
  \draw[arr] (s) -- (r);  \draw[arr] (r) -- (o);
\end{tikzpicture}
\caption{EHR Analyst — the analyse/structure/review three-call sequence
applied within the generic agent blueprint. Each call is a separate LLM
invocation with its own system prompt loaded from
\texttt{prompts/ehr\_analyst.yaml}.}
\label{fig:agent_ehr_flow}
\end{figure}

\subsection{Lab Interpreter}
\label{sec:agent_lab}
The Lab Interpreter consumes \texttt{lab\_case.json} and classifies each
laboratory observation against reference ranges, interprets temporal trends,
and ranks findings by clinical severity. Critical values are flagged with an
explicit \texttt{CRITICAL} marker that propagates downstream. Like the EHR
Analyst, the Lab Interpreter runs the three-call analyse--structure--review
sequence (Figure~\ref{fig:agent_lab_flow}); the structure call additionally
applies a deterministic severity-ranking rule that orders findings by their
deviation from reference range, putting critical alerts first.

\begin{figure}[ht]
\centering
\begin{tikzpicture}[
  every node/.style={font=\scriptsize},
  inp/.style={rectangle, draw, rounded corners=2pt, minimum height=10mm,
              minimum width=24mm, align=center, fill=yellow!30},
  llm/.style={rectangle, draw, rounded corners=2pt, minimum height=10mm,
              minimum width=24mm, align=center, fill=blue!12},
  rank/.style={rectangle, draw, rounded corners=2pt, minimum height=10mm,
               minimum width=22mm, align=center, fill=orange!18,
               font=\scriptsize\itshape},
  outbox/.style={rectangle, draw, rounded corners=2pt, minimum height=10mm,
              minimum width=28mm, align=center, fill=green!18},
  arr/.style={-{Latex[length=1.6mm]}, semithick}
]
  \node[inp] (in) {\texttt{lab\_case.json}\\(Gold object)};
  \node[llm, right=4mm of in] (a)
    {Call 1: \textbf{analyse}\\\scriptsize classify, trend};
  \node[llm, right=4mm of a] (s)
    {Call 2: \textbf{structure}\\\scriptsize JSON + flags};
  \node[rank, below=4mm of s] (rule)
    {severity rule\\(deterministic)};
  \node[llm, right=4mm of s] (r)
    {Call 3: \textbf{review}\\\scriptsize verify ranking};
  \node[outbox, right=4mm of r] (o)
    {\texttt{lab\_interpreter}\\\scriptsize ranked findings\\\scriptsize CRITICAL flags};
  \draw[arr] (in) -- (a); \draw[arr] (a) -- (s);
  \draw[arr] (s) -- (r);  \draw[arr] (r) -- (o);
  \draw[arr, dashed] (s) -- (rule);
  \draw[arr, dashed] (rule.east) -| (r.south);
\end{tikzpicture}
\caption{Lab Interpreter — same three-call blueprint as the EHR Analyst,
with an additional deterministic severity-ranking step injected between
the structure and review calls. The ranking rule orders findings by their
deviation from reference range and propagates a \texttt{CRITICAL} marker
to downstream agents.}
\label{fig:agent_lab_flow}
\end{figure}

\subsection{Diagnostic Reasoning}
\label{sec:diagnostic_reasoning}
\label{sec:agent_diag}
The Diagnostic Reasoning agent is the central hypothesis-generating
component of the pipeline. Unlike Stage~1 agents, which execute a fixed
three-call sequence, this agent implements two distinct reasoning modes
selected at run time: an \textit{adaptive critique loop} when the
multi-level memory subsystem is disabled or no quality priors are
available, and a \textit{three-phase split} when Tier-4 case-based
recall returns at least one DIRECT-match prior above the configured
similarity threshold (Section~\ref{sec:tier4}). Both modes are bounded
to at most three reasoning rounds.

The adaptive critique loop proceeds as follows. A first call synthesises
evidence from the Stage~1 outputs, a second call generates at least
three candidate diagnoses, a third call ranks them by probability, and
subsequent critique rounds re-rank the differential until either a
confidence threshold of 75 is reached or three rounds have elapsed.
This adaptive behaviour was motivated by the finding of Shang et
al.~\cite{Shang2025DynamiCare} that dynamic multi-round reasoning
significantly outperforms single-pass approaches on complex cases, and by
the MDAgents architectural precedent of Kim et al.~\cite{Kim2024MDAgents}
that complexity-based routing improves both accuracy and efficiency.

\paragraph{Three-phase split.}
When the memory subsystem returns at least one quality prior, the
adaptive critique loop is replaced by three explicit re-ranking rounds
designed to harvest the prior without anchoring on it
(Figure~\ref{fig:agent_diag_flow}).  \textit{Round~A} re-ranks the
differential with the case-based prior visible and the agent
instructed to weight the prior but anchor every diagnosis to evidence
in the current patient's data. \textit{Round~B} re-ranks
\emph{clean-room}, with the prior suppressed and an explicit
anti-anchoring instruction (``reason from the patient's evidence
alone, as if this were the first patient you have ever seen'').
\textit{Round~C} synthesises rounds A and B: where they agree, the
consensus is emitted; where they disagree, the agent prefers Round~B
unless Round~A's choice has a clearly named supporting finding in the
current case. The split is informed by the same multi-round literature
above and by the \textit{silent agreement} failure mode identified by
Wang et al.~\cite{Wang2025Catfish}; the explicit clean-room round is
this work's contribution.

\begin{figure}[ht]
\centering
\begin{tikzpicture}[
  every node/.style={font=\scriptsize},
  inp/.style={rectangle, draw, rounded corners=2pt, minimum height=8mm,
              minimum width=22mm, align=center, fill=yellow!30},
  fixed/.style={rectangle, draw, rounded corners=2pt, minimum height=8mm,
                minimum width=24mm, align=center, fill=blue!12},
  rA/.style={rectangle, draw, rounded corners=2pt, minimum height=14mm,
             minimum width=30mm, align=center, fill=violet!18},
  rB/.style={rectangle, draw, rounded corners=2pt, minimum height=14mm,
             minimum width=30mm, align=center, fill=blue!18},
  rC/.style={rectangle, draw, rounded corners=2pt, minimum height=14mm,
             minimum width=30mm, align=center, fill=teal!18},
  prior/.style={rectangle, draw, dashed, rounded corners=2pt,
                minimum height=8mm, minimum width=30mm, align=center,
                fill=violet!8, font=\scriptsize\itshape},
  outbox/.style={rectangle, draw, rounded corners=2pt, minimum height=8mm,
              minimum width=24mm, align=center, fill=green!18},
  alt/.style={rectangle, draw, rounded corners=2pt, minimum height=14mm,
              minimum width=30mm, align=center, fill=blue!8,
              font=\scriptsize\itshape},
  arr/.style={-{Latex[length=1.6mm]}, semithick},
  loop/.style={-{Latex[length=1.4mm]}, dashed, semithick, blue!60!black},
  branch/.style={font=\scriptsize\bfseries, text=violet!60!black}
]
  % Top row: fixed pre-rounds
  \node[inp] (s1) {EHR + Lab\\summaries};
  \node[fixed, right=3mm of s1] (c1)
    {Call 1: \textbf{synthesis}};
  \node[fixed, right=3mm of c1] (c2)
    {Call 2: \textbf{hypotheses}};
  \node[fixed, right=3mm of c2] (c3)
    {Call 3: \textbf{initial rank}};
  \draw[arr] (s1) -- (c1); \draw[arr] (c1) -- (c2); \draw[arr] (c2) -- (c3);

  % Branch decision
  \node[branch, below=4mm of c3, align=center] (decision)
    {if Tier-4 quality priors\\\textnormal{\scriptsize (DIRECT, score $\geq 0.75$)}};
  \draw[arr] (c3) -- (decision);

  % LEFT branch: 3-phase split (when priors exist)
  \node[prior, below left=8mm and -8mm of decision] (prior)
    {Tier-4 prior block};
  \node[rA, below=4mm of prior] (rA)
    {\textbf{Round A}\\WITH prior\\\scriptsize re-rank};
  \node[rB, below=4mm of rA] (rB)
    {\textbf{Round B}\\CLEAN-ROOM\\\scriptsize anti-anchoring};
  \node[rC, below=4mm of rB] (rC)
    {\textbf{Round C}\\synthesise A + B};
  \draw[arr, violet!60!black] (decision.west) -- ++(-12mm,0) |- (prior.east);
  \draw[arr] (prior) -- (rA);
  \draw[arr] (rA) -- (rB);
  \draw[arr] (rB) -- (rC);

  % RIGHT branch: adaptive critique loop (fallback)
  \node[alt, below right=8mm and -8mm of decision] (alt)
    {\textbf{adaptive critique loop}\\$\leq 3$ rounds\\\scriptsize stop at conf $\geq 75$};
  \draw[arr] (decision.east) -- ++(12mm,0) |- (alt.east);
  \draw[loop] ($(alt.north east)+(-3mm,-1mm)$)
    .. controls +(right:8mm) and +(right:8mm) ..
    ($(alt.south east)+(-3mm,1mm)$);

  % Final structured-JSON call (always)
  \node[fixed, below=12mm of decision, xshift=0mm, fill=red!12] (jcall)
    {Final call: \textbf{structured JSON}};
  \draw[arr] (rC.east) -| (jcall.west);
  \draw[arr] (alt.west) -| (jcall.east);

  \node[outbox, below=4mm of jcall] (o)
    {\texttt{diagnostic\_reasoning}\\\scriptsize ranked differential};
  \draw[arr] (jcall) -- (o);

\end{tikzpicture}
\caption{Diagnostic Reasoning agent — fixed three-call pre-rounds
(synthesis, hypothesis, initial rank), then a runtime branch based on
the availability of Tier-4 quality priors. The 3-phase split (left)
runs when at least one DIRECT-match prior with cosine similarity
$\geq 0.75$ is recalled; the adaptive critique loop (right) runs
otherwise. Both paths terminate in a structured-JSON final call that
emits the agent's contribution to \texttt{agent\_outputs}.}
\label{fig:agent_diag_flow}
\end{figure}

\subsection{Clinical Reviewer}
\label{sec:agent_reviewer}
The Clinical Reviewer is an adversarial agent that independently re-examines
the evidence and the diagnostic agent's differential. It is explicitly
prompted to challenge the primary diagnosis, surface plausible alternatives
that may have been missed, and identify cognitive biases in the diagnostic
agent's reasoning. This design directly addresses the \textit{silent
agreement} failure mode identified by Wang et al.\ in the Catfish Agent
framework~\cite{Wang2025Catfish}: by instantiating the reviewer with an
adversarial rather than confirmatory system prompt, the pipeline forces
genuine deliberation rather than reflexive consensus.

The agent's internal flow is shown in Figure~\ref{fig:agent_reviewer_flow}.
With the multi-level memory subsystem enabled, the Reviewer additionally
reads the Diagnostic agent's Tier-2 critique trail and uses it during
the verification call to challenge the \emph{path} to the diagnosis,
not only the diagnosis itself.

\begin{figure}[ht]
\centering
\begin{tikzpicture}[
  every node/.style={font=\scriptsize},
  inp/.style={rectangle, draw, rounded corners=2pt, minimum height=10mm,
              minimum width=26mm, align=center, fill=yellow!30},
  llm/.style={rectangle, draw, rounded corners=2pt, minimum height=10mm,
              minimum width=24mm, align=center, fill=blue!12},
  mem/.style={rectangle, draw, dashed, rounded corners=2pt, minimum height=8mm,
              minimum width=26mm, align=center, fill=blue!8,
              font=\scriptsize\itshape},
  outbox/.style={rectangle, draw, rounded corners=2pt, minimum height=10mm,
              minimum width=28mm, align=center, fill=green!18},
  arr/.style={-{Latex[length=1.6mm]}, semithick},
  memarr/.style={-{Latex[length=1.4mm]}, dashed, blue!60!black, thin}
]
  \node[inp] (in) {Stage 1 + 2\\agent outputs};
  \node[llm, right=4mm of in] (a)
    {Call 1: \textbf{re-analyse}\\\scriptsize independent read};
  \node[llm, right=4mm of a] (v)
    {Call 2: \textbf{verify}\\\scriptsize per-diagnosis};
  \node[llm, right=4mm of v] (j)
    {Call 3: \textbf{JSON}\\\scriptsize verdict + flags};
  \node[outbox, right=4mm of j] (o)
    {\texttt{clinical\_reviewer}\\\scriptsize CONFIRM\,/\,MODIFY\,/\,REJECT};
  \draw[arr] (in) -- (a); \draw[arr] (a) -- (v);
  \draw[arr] (v) -- (j);  \draw[arr] (j) -- (o);

  % Tier-2 read into the verification call
  \node[mem, above=5mm of v] (t2)
    {Tier-2: Diagnostic\\critique trail};
  \draw[memarr] (t2) -- (v);
\end{tikzpicture}
\caption{Clinical Reviewer — three-call sequence (independent re-analysis,
per-diagnosis verification, structured-JSON output) augmented by a Tier-2
read of the Diagnostic agent's critique trail (dashed) so the verification
call can challenge the reasoning path, not only the final differential.}
\label{fig:agent_reviewer_flow}
\end{figure}

\subsection{Diagnostic Refiner}
\label{sec:agent_refiner}
The Diagnostic Refiner merges the Diagnostic Reasoning and Clinical Reviewer
outputs into a single final differential, reconciling any disagreement between
the two. When the two agents agree, the refiner acts as a light rewriter that
harmonises terminology and adjusts rank positions. When they disagree, the
refiner applies a small decision procedure: per-diagnosis reviewer flags of
\textit{verified} promote the diagnosis, \textit{refuted} demote it, and
\textit{uncertain} retain its position. The final output is written to
\texttt{agent\_outputs.final\_diagnosis}. Figure~\ref{fig:agent_refiner_flow}
sketches the merge step; with multi-level memory enabled, the refiner
additionally reads the Tier-2 session timeline and any Tier-3 priors for
the diseases appearing in the candidate differential.

\begin{figure}[ht]
\centering
\begin{tikzpicture}[
  every node/.style={font=\scriptsize},
  inp/.style={rectangle, draw, rounded corners=2pt, minimum height=10mm,
              minimum width=28mm, align=center, fill=yellow!30},
  llm/.style={rectangle, draw, rounded corners=2pt, minimum height=10mm,
              minimum width=30mm, align=center, fill=blue!12},
  mem/.style={rectangle, draw, dashed, rounded corners=2pt, minimum height=8mm,
              minimum width=30mm, align=center, fill=teal!12,
              font=\scriptsize\itshape},
  outbox/.style={rectangle, draw, rounded corners=2pt, minimum height=10mm,
              minimum width=28mm, align=center, fill=green!18},
  arr/.style={-{Latex[length=1.6mm]}, semithick},
  memarr/.style={-{Latex[length=1.4mm]}, dashed, teal!60!black, thin}
]
  \node[inp] (diag) {Diagnostic Reasoning\\differential};
  \node[inp, below=3mm of diag] (rev) {Clinical Reviewer\\verdict + flags};
  \node[llm, right=6mm of $(diag)!0.5!(rev)$] (merge)
    {\textbf{merge call}\\\scriptsize verified $\uparrow$, refuted $\downarrow$,\\\scriptsize uncertain held};
  \node[outbox, right=6mm of merge] (out) {\texttt{final\_diagnosis}};
  \draw[arr] (diag.east) -- (merge.west);
  \draw[arr] (rev.east)  -- (merge.west);
  \draw[arr] (merge) -- (out);

  % Memory reads
  \node[mem, above=4mm of merge] (t2)
    {Tier-2 timeline};
  \node[mem, below=4mm of merge] (t3)
    {Tier-3 priors\\(per disease in differential)};
  \draw[memarr] (t2) -- (merge);
  \draw[memarr] (t3) -- (merge);
\end{tikzpicture}
\caption{Diagnostic Refiner — single merge call producing the final
differential. With multi-level memory enabled, it additionally reads
the Tier-2 session timeline (above) and any Tier-3 disease-level priors
for the candidate diagnoses (below); both are dashed because they are
optional inputs depending on the \texttt{MEMORY\_ENABLED} flag.}
\label{fig:agent_refiner_flow}
\end{figure}

\subsection{LLM Evaluator}
\label{sec:agent_eval}
The LLM Evaluator compares the final differential against the Synthea ground
truth stored in \texttt{ground\_truth.json}. The evaluator is deliberately
invoked on a separate LLM configuration (\texttt{LLM\_EVALUATOR\_MODEL}) from
the reasoning agents, following the methodological convention that the judge
should not also be the examinee. The evaluator uses the LLM-as-Judge
methodology described in Section~\ref{sec:evaluation} and assigns one of three
match classifications: DIRECT (semantic equivalent of the target), INDIRECT
(clinically related precursor, consequence, or subtype), or MISS (unrelated).

\subsection{Treatment Planning}
\label{sec:agent_tx}
The Treatment Planning agent runs only when the evaluator has confirmed a
DIRECT match, ensuring that guideline-grounded treatment recommendations are
produced only for correctly identified conditions. It uses Retrieval-Augmented
Generation over a Qdrant vector store of NICE clinical guidelines
(Section~\ref{sec:rag}) to retrieve the top-$k$ relevant guideline passages,
which are then passed to the LLM as context alongside the patient's final
diagnosis. The agent emits a structured plan containing first-line
pharmacotherapy, non-pharmacological measures, a monitoring schedule, and
red-flag symptoms.

% ─────────────────────────────────────────────────────────────────────────────
\section{Retrieval-Augmented Generation Layer}
\label{sec:rag}
% ─────────────────────────────────────────────────────────────────────────────

The Treatment Planning agent is the only agent in the pipeline that consults
external knowledge at inference time. This deliberate isolation of retrieval
to the treatment stage keeps the diagnostic reasoning grounded exclusively in
the patient-specific evidence, and defers external guideline material to the
point at which it is most clinically appropriate: after a diagnosis has been
established and verified.

The retrieval corpus consists of NICE (National Institute for Health and Care
Excellence) clinical guideline documents relevant to the eight target disease
categories. Guidelines are chunked into overlapping passages of approximately
400 tokens with 50-token overlap, embedded with the \textit{BioLORD-2023}
biomedical sentence encoder, and stored in a Qdrant vector collection indexed
by cosine similarity. The choice of Qdrant over simpler alternatives such as
in-memory cosine search was motivated by its support for filtered retrieval:
the treatment agent restricts retrieval to passages tagged with the final
diagnosis's ICD-10 or SNOMED code, which materially reduces the risk of
guideline hallucination.

At inference time, the treatment agent issues a query composed of the final
diagnosis name and the top three supporting findings. The top-five passages
by cosine similarity are retrieved, deduplicated, and inserted into the LLM
prompt along with an explicit instruction to cite the retrieved passage
identifier for every concrete recommendation. The Output Gate rejects any
plan that references medications or procedures not present in the retrieved
passages, providing a second layer of defence against hallucinated
recommendations.

% ─────────────────────────────────────────────────────────────────────────────
\section{LLM Configuration and Provider Abstraction}
\label{sec:llm_config}
% ─────────────────────────────────────────────────────────────────────────────

To permit experimentation with multiple language models without code changes,
all LLM calls are routed through a single provider-agnostic adapter that
implements a uniform \texttt{get\_llm()} interface. The adapter exposes five
backend providers --- Groq, OpenAI, Anthropic, Google Gemini, and Ollama ---
behind a common factory that reads configuration from environment variables.
A second factory, \texttt{get\_evaluator\_llm()}, returns a separately
configured model for the Evaluator, which allows the judge to run on a
different provider from the reasoning agents when required by the
experimental protocol.

Three model configurations were used in this thesis:
\begin{itemize}
  \item \textbf{GPT-OSS 120B} served via the Groq cloud API, used as the
  primary reasoning model.
  \item \textbf{Med42 70B} (\texttt{thewindmom/llama3-med42-70b}), a Llama-3
  checkpoint fine-tuned on clinical text, served locally via Ollama and used
  in the medical-fine-tuning comparison reported in
  Chapter~\ref{chap:results}.
  \item \textbf{Qwen3 32B} served via Groq, used exclusively as the LLM
  Evaluator judge. Qwen3 32B was chosen because it is independent of the two
  reasoning models under test, preventing the judging result from being
  biased by shared training data.
\end{itemize}

All LLM calls are subject to an exponential back-off retry policy capped at
three attempts and a per-call timeout of 120 seconds. The LangSmith tracing
integration is enabled for every call, producing an audit trail of prompts,
responses, and intermediate errors that is stored in a cloud-hosted
observability platform.

% ─────────────────────────────────────────────────────────────────────────────
\section{Evaluation Framework}
\label{sec:evaluation}
% ─────────────────────────────────────────────────────────────────────────────

Evaluation of the pipeline is driven by two design principles. First, the
ground truth must be independent of any LLM: because the pipeline itself is
composed of LLMs, using an LLM-generated ground truth would produce a
circular evaluation that is biased toward the peculiarities of the particular
generator model. Synthea's rule-based state machines provide this
independence. Second, the evaluator should not be the examinee: using the
same model to generate and judge produces inflated self-consistency scores
that overestimate real performance. The framework therefore enforces
separation between the reasoning model and the judge model.

\subsection{Ground Truth Extraction}
For every Synthea patient, a \texttt{ground\_truth.json} file is materialised
during Gold assembly. It contains the single \textit{target condition} that
the patient was intended to carry, expressed as both a SNOMED concept and a
human-readable string, together with the date of first diagnosis. The file
is read only by the LLM Evaluator in Stage~5 and is never exposed to any
reasoning agent.

\subsection{LLM-as-Judge Methodology}
The LLM Evaluator applies a structured judging prompt that presents the
target disease together with the top-five diagnoses from the pipeline's final
differential. The judge classifies each candidate as DIRECT, INDIRECT, or
UNRELATED according to a fixed rubric with explicit worked examples, then
selects the best-matching candidate and reports its rank. Temperature is
fixed at zero and the response schema is a five-line structured format
parsed by a regex-based parser. Table~\ref{tab:match_types} summarises the
rubric.

\begin{table}[ht]
\centering
\caption{Match-type rubric applied by the LLM Evaluator.}
\label{tab:match_types}
\small
\begin{tabular}{@{}p{22mm}p{95mm}@{}}
\toprule
\textbf{Match type} & \textbf{Criterion} \\
\midrule
DIRECT   & The candidate names the same disease as the ground truth, possibly
           under a clinically equivalent synonym (e.g.\ ``Atherosclerotic CVD''
           is a DIRECT match for ``Ischemic heart disease''). \\
INDIRECT & The candidate names a cause, consequence, precursor, or subtype of
           the ground truth (e.g.\ ``CKD stage 4'' is INDIRECT for
           ``End-stage renal disease''). \\
MISS     & The candidate has no clinical connection to the ground truth. \\
\bottomrule
\end{tabular}
\end{table}

\subsection{Metrics}
The primary metrics reported in Chapter~\ref{chap:results} are:
\begin{itemize}
  \item \textbf{Direct rate} --- the proportion of patients for which any
  diagnosis in the top five is classified as DIRECT.
  \item \textbf{Found rate} --- combined DIRECT plus INDIRECT rate.
  \item \textbf{Miss rate} --- the proportion of patients with no matching
  diagnosis in the top five.
  \item \textbf{Top-1 direct rate} --- the proportion for which the primary
  (rank-1) diagnosis is classified as DIRECT. This is the strictest metric
  and best reflects clinical usefulness.
  \item \textbf{Agent success rate} --- per-agent, the proportion of patients
  for which the agent terminated without error.
  \item \textbf{Per-patient latency} --- wall-clock time from orchestrator
  invocation to final-diagnosis write.
\end{itemize}

% ─────────────────────────────────────────────────────────────────────────────
\section{Hardware and Software Environment}
\label{sec:env}
% ─────────────────────────────────────────────────────────────────────────────

Experiments were executed on a local Apple Silicon workstation for the data
pipeline and for Ollama-hosted local models (Med42 70B), and on the Groq
cloud inference platform for the Groq-hosted models (GPT-OSS 120B and
Qwen3 32B). The Qdrant vector store is hosted on a managed cloud deployment
to avoid local persistence overhead. Table~\ref{tab:software_stack} lists the
principal software dependencies.

\begin{table}[ht]
\centering
\caption{Principal software dependencies.}
\label{tab:software_stack}
\small
\begin{tabular}{@{}p{34mm}p{34mm}p{58mm}@{}}
\toprule
\textbf{Component} & \textbf{Technology} & \textbf{Role} \\
\midrule
Orchestration   & LangGraph              & Typed \texttt{StateGraph} with checkpointing \\
Agent runtime   & LangChain (LCEL)       & Prompt composition and LLM invocation \\
LLM providers   & Groq, Ollama           & Cloud and local inference back-ends \\
Vector store    & Qdrant                 & NICE guideline retrieval \\
Embedding model & BioLORD-2023           & Biomedical sentence encoder \\
Output schemas  & Pydantic v2            & Structured agent outputs \\
Data store      & DuckDB                 & Local OMOP CDM warehouse \\
Transformation  & dbt-core               & SQL transformations with tests \\
Ingestion       & Python + PyArrow       & Parquet ingestion and JSON assembly \\
DAG runner      & Prefect                & Data pipeline orchestration \\
Logging         & structlog              & JSON-lines structured logs \\
Observability   & LangSmith              & LLM call tracing \\
Portal          & Streamlit              & Read-only review dashboard \\
\bottomrule
\end{tabular}
\end{table}

All environment-specific configuration (provider, model name, API keys,
Qdrant endpoint, output directory) is externalised into a single
\texttt{.env} file read by a central \texttt{Config} class at process start.
This design ensures that experimental variations --- for example, switching
from GPT-OSS to Med42 --- are a configuration change rather than a code
change, supporting reproducibility of individual runs.

% ─────────────────────────────────────────────────────────────────────────────
\section{End-to-End Execution Pipeline}
\label{sec:pipeline_e2e}
% ─────────────────────────────────────────────────────────────────────────────

The complete execution path of a single patient is summarised in
Algorithm~\ref{alg:e2e}. The pipeline is invoked either through the command
line (for single-patient or batch processing) or through the Streamlit
portal (for interactive case review).

\begin{algorithm}[ht]
\caption{End-to-end processing of a single patient.}
\label{alg:e2e}
\begin{algorithmic}[1]
\STATE Load Gold case objects \texttt{ehr\_case.json} and \texttt{lab\_case.json}
\STATE Initialise \texttt{PipelineState} with patient context and empty output slots
\STATE \textbf{Stage 1 (parallel):} invoke EHR Analyst and Lab Interpreter concurrently
\STATE \textbf{Stage 2:} invoke Diagnostic Reasoning; repeat critique round until
       confidence $\geq 75$ or three rounds executed
\STATE \textbf{Stage 3:} invoke Clinical Reviewer with Stage 1 and Stage 2 outputs
\STATE \textbf{Stage 4:} invoke Diagnostic Refiner to merge Stages 2 and 3
\STATE \textbf{Stage 5:} invoke LLM Evaluator to classify the final diagnosis
       against the Synthea ground truth
\IF{evaluation match type $=$ DIRECT}
  \STATE \textbf{Stage 6:} invoke Treatment Planning with RAG over Qdrant
\ELSE
  \STATE Skip Stage 6 (treatment planning deferred until correct diagnosis)
\ENDIF
\STATE Persist \texttt{agent\_outputs}, \texttt{execution\_trace}, and per-agent
       JSON files to \texttt{data/gold/mas\_results/<patient\_uuid>/}
\end{algorithmic}
\end{algorithm}

All agent outputs, the execution trace, and the evaluation result are
persisted to disk after each patient, providing a durable record that can
be reloaded by the Streamlit portal or re-analysed offline without re-running
the pipeline.
